% ============================================================
%  Control Reflectors Scan  —  Testing and Results
% ============================================================

\subsection{Control Reflectors Scan}

Before acquiring any foliage scans, a free-space baseline measurement was performed on
11~March~2026 under clear conditions with no vegetation in the radar beam path.
The purpose of this scan was twofold: to verify that both corner reflectors (CRs)
produced strong, well-separated returns under unobstructed conditions, and to characterise
the ambient clutter floor and free-space propagation behaviour of the radar system.
These reference values are used throughout the subsequent analysis to compute
vegetation attenuation and to define detection thresholds for the foliage scans.

% ---- LiDAR ----
\subsubsection*{LiDAR Identification of Corner Reflector Positions}

The LiDAR scan of the control scene confirmed the physical positions of both reflectors.
Raw point cloud returns near the reflectors were dominated by grass and ground clutter,
so a height filter was applied to isolate high-intensity returns above the ground plane.
Figure~\ref{fig:lidar_control_identified} shows the resulting point cloud
after filtering, with Reflector~1 and Reflector~2 highlighted.
The centroid of each reflector cluster was computed as its mean Cartesian position:

\begin{itemize}
  \item Reflector~1: $(8.71,\ {-1.95},\ 0.33)$~m
  \item Reflector~2: $(9.66,\ 2.77,\ 0.42)$~m
\end{itemize}

The LiDAR-derived separation between the two reflectors was 4.82~m.
These positions served as the ground-truth reference for spatial co-registration
with the radar data.

% ---- Radar ----
\subsubsection*{Radar Control Scan}

Figure~\ref{fig:control_overview} summarises the full radar control scan analysis
across three panels.

\paragraph{Amplitude vs.\ range.}
The upper panel plots every recorded return as a function of range.
To visualise the overall trend without being dominated by individual outliers,
returns were grouped into 2~m range bins and the bin mean $\pm 1\sigma$ is
overlaid as black data points with error bars.
Both CRs appear as narrow, high-amplitude clusters that stand unambiguously above
the surrounding clutter: CR1 produces a peak of 66~dB at a range of approximately
19.8~m (shaded red), and CR2 produces a peak of 34~dB at approximately 22~m
(shaded green).
The clutter floor, defined as the mean amplitude of all background returns
excluding the CR clusters, was $-27.8$~dB and is shown as the dashed orange line.
CR1 therefore sits 93.8~dB above the clutter floor, and CR2 sits 61.8~dB above it.
The 32~dB difference between the two peak amplitudes reflects the difference in
radar cross-section between the large and small reflectors, compounded by their
slightly different ranges.

A linear fit to all returns yielded a free-space propagation slope of
$-0.303 \pm 0.308$~dB/m ($R^2 = 0.05$, $p = 0.337$).
The near-zero $R^2$ and non-significant $p$-value confirm that amplitude does not
vary systematically with range in the absence of vegetation, as expected for a
flat clutter background. This validates $-27.8$~dB as a stable, range-independent
reference for all subsequent foliage scans.

The range separation between the two CRs measured by the radar was 3.98~m,
compared with 4.82~m derived from LiDAR. The sub-metre discrepancy arises because
the radar range corresponds to the peak-scattering point on each reflector, whereas
the LiDAR centroid is computed over the full illuminated reflector surface.

\paragraph{Amplitude distribution.}
The lower-left panel shows a histogram of all return amplitudes.
The bulk of returns form a roughly symmetric distribution centred near $-25$~dB,
consistent with the estimated clutter floor of $-27.8$~dB.  The two CR peaks
(66~dB for CR1, 34~dB for CR2, marked by dotted vertical lines) are clear outliers
at the high end of the distribution, separated from the clutter body by a gap of
more than 60~dB for CR1 and approximately 60~dB for CR2.
The well-separated CR peaks confirm that a simple amplitude threshold is sufficient
to detect both reflectors with no false positives in free-space conditions.

\paragraph{XY point cloud.}
The lower-right panel shows all returns projected onto the horizontal plane and
coloured by amplitude.  High-amplitude returns (yellow, above roughly $+20$~dB)
are tightly clustered at the known positions of CR1 and CR2, marked by green stars.
The remainder of the scene is occupied by low-amplitude clutter (purple, below
$-20$~dB), confirming that the strong returns are spatially localised to the
reflector locations and are not artefacts distributed across the scene.

% ---- Figure caption (place near \includegraphics in main file) ----
% \begin{figure}[htbp]
%   \centering
%   \includegraphics[width=\linewidth]{figures/control_scan_overview}
%   \caption{Control Scan --- Free Space Baseline (11~March~2026, no foliage).
%     \textit{Top:} Amplitude versus range for all returns.  CR1 (66~dB, red shaded
%     region) and CR2 (34~dB, green shaded region) are visible as sharp peaks above
%     the clutter floor of $-27.8$~dB (dashed orange line).  Bin means $\pm 1\sigma$
%     are shown as black data points.  The linear fit (dashed black line,
%     $-0.303$~dB/m, $R^2 = 0.05$) confirms a range-independent clutter background.
%     \textit{Bottom left:} Amplitude distribution of all returns; CR peaks appear as
%     clear outliers above the clutter body.
%     \textit{Bottom right:} XY point cloud coloured by amplitude; high-amplitude
%     returns are spatially concentrated at the CR locations (green stars).}
%   \label{fig:control_overview}
% \end{figure}

% ---- Summary ----
\subsubsection*{Baseline Summary}

The control scan established the following reference values used throughout all
subsequent analysis:

\begin{itemize}
  \item \textbf{Clutter floor:} $-27.8$~dB (mean amplitude of background returns
        in the absence of vegetation)
  \item \textbf{CR1 peak amplitude:} $66$~dB at 19.8~m range (SNR $= 93.8$~dB)
  \item \textbf{CR2 peak amplitude:} $34$~dB at 23.8~m range (SNR $= 61.8$~dB)
  \item \textbf{Free-space propagation slope:} $-0.303$~dB/m (not statistically
        significant; amplitude is range-independent in free space)
\end{itemize}

The large SNR margins — 93.8~dB for CR1 and 61.8~dB for CR2 — demonstrate that
both reflectors are easily detectable in unobstructed conditions.
Any reduction in these margins in the foliage scans can therefore be attributed
directly to vegetation attenuation rather than to system noise or sensor limitations.
